% ============ 6. Robustness and Threats to Validity ============
\section{Robustness and Threats to Validity}\label{sec:robust}

\subsection{Specification robustness of the baseline}

Five checks, all computed fresh on the archived data. (i)~Raising the relevance threshold from 30 to 50 leaves the baseline essentially unchanged ($t=2.12$, $N=187$). (ii)~Using the unweighted mean over all headlines with no relevance filter gives $t=2.07$ ($N=209$): the filter is not doing the work. (iii)~Dropping dense event days (five or more covered firms) leaves the point estimate intact but the statistic at $t=1.61$ ($N=117$): about a quarter of the baseline's strength comes from clustering on event days, consistent with H5. (iv)~The anonymized condition with the same-day control runs at $t=1.59$, tracking its named twin. (v)~A within-day permutation placebo (1{,}000 draws, sentiment permuted inside each trading day) yields a two-sided permutation $p=0.007$: whatever the baseline coefficient is, it is not an artifact of arbitrary re-labeling. The consistent picture across (i)--(v) is a real but fragile association~---robust to measurement choices, weakened by event-day down-weighting, and significant at conventional levels only under the original sample construction.

\subsection{Session robustness of the predictive coefficient}

The baseline coefficient rests on October-15 scores; the three October-16 named repeats re-score the identical headlines in a different sitting, the October-18 interleaved session re-scores them twice more per condition in a third, and the session-4 mechanism passes add three further named draws at both batch compositions plus a random re-batch. Across the fourteen archived named draws~---twelve individual passes (the original, three October-16 repeats, four October-18 rounds, the October-18 542-composition pass, and the three session-4 passes) plus the two session means~---the next-day coefficient lives in $[0.010, 0.016]$ and the $t$-statistic in $[1.47, 2.12]$: the point estimate is stable across sessions, compositions, and re-batching to well within its own standard error, and every repeat falls inside the range the specification checks above already span~---but the \emph{significance} of the coefficient moves with the draw, from $p=0.035$ (the baseline round) to $p=0.14$ (the weakest October-18 round). The same-day loading is large under every draw ($t$ between $2.96$ and $3.80$ for the session-1/2 draws). Session drift and batch composition are therefore \emph{level} phenomena~---they move the mean score by points (Section~\ref{sec:results:h2})~---not \emph{slope} phenomena: predictive loadings survive re-scoring, which is the robustness fact every other result in this paper leans on, and the composition contrasts of the slope are individually insignificant ($p>0.9$; Appendix~\ref{app:interleaved}). The same-session instrument difference is itself stable across all three designs: the anonymized-versus-named predictive difference is $0.0002$ ($p=0.905$) inside the October-16 sitting, $0.0013$ ($p=0.622$) inside the interleaved October-18 session, tracking its cross-session twin ($p=0.887$).

\subsection{Auxiliary analyses}
\label{sec:robust:aux}

Three descriptive exercises that a predictability paper would normally foreground~---the event-horizon path of the coefficient, sentiment-sorted quintile portfolios, and heterogeneity interactions~---are relegated to Appendix~\ref{sec:appaux} and summarized here in one sentence each, because none survives the audit's inference standard as an independent finding. The horizon path (positive at $t+1$, negative at $t+2$ with $t=-2.04$, fading by $t+3$) is mechanically consistent with continuation-and-reversal around high-absolute-return events on a sample selected on exactly those events \citep{chan2003}, and we read it as the echo channel propagating through return dynamics rather than as evidence of attention-driven overshooting. The quintile spread (Q5$-$Q1 of 89.9 \bps{}) is non-monotone across intermediate quintiles and statistically insignificant (Welch $t=1.24$, $p=0.22$), with null information-coefficient and direction tests. Of the three heterogeneity interactions, only the firm-size interaction is significant ($t=2.25$); the news-intensity interaction is marginal ($t=1.90$) and the volatility interaction is indistinguishable from zero ($t=0.01$)~---the split-sample rhetoric that would have read the last of these as a ``high-volatility concentration'' is exactly the practice Section~\ref{sec:framework:inference} rejects.

\subsection{Threats to validity we could not address}
\label{sec:robust:threats}

Five threats remain outside the audit's reach. \emph{Publication-timing granularity}: headlines carry dates but not timestamps, so the same-day return attributed to a news day may include pre-publication price movement, and we cannot separate before-open, intraday, and after-close news~---the overnight/intraday split is a coarse proxy, and all timing claims built on it are exploratory. \emph{Single scorer, and single served model}: all conditions use one model, \texttt{glm-4-plus}; cross-model stability (the model-variability problem of \citealp{herrera2025}) is not tested here. We record a further constraint discovered while preparing the control sessions: the serving API ignores the client-side model identifier~---every requested model string is served by the same \texttt{glm-4-plus} endpoint (the probe is archived in the replication package)~---so a cross-model contrast was not executable in this environment at all, and the within-model estimates bound nothing about other models. The audit rounds ran twelve hours after the original round, and the identical named protocol drifted $+5.8$ points in level between sittings; the interleaved session has since decomposed that drift into a $+3.6$-point batch-composition component and a non-monotonic endpoint component of about $\pm1$ point per session-pair (Section~\ref{sec:results:h2}), so the drift is no longer unattributed~---but its endpoint component remains outside our control, and the repeat-call ICC therefore bounds near-term stability only. \emph{No human gold standard}: without human sentiment labels we can audit the score's relation to returns and to its own perturbations, but not its semantic accuracy against ground truth. The price-language classification audit of Section~\ref{sec:results:h4} is LLM-assisted and does not substitute for one; a complete two-rater human annotation package (codebook, per-item sheets, and an inter-rater reliability script) is prepared and archived in the replication package so that human labels can be added without re-design, but no human labels existed when this paper was written. \emph{Alignment.} Publication dates without timestamps prevent event-window precision; same-day returns conflate pre- and post-publication movement; the overnight/intraday decomposition is coarse. Mechanism claims about timing (attention diffusion versus liquidity) are therefore not identifiable in this design, which is why the paper makes none. \emph{External validity}: 36 mega-caps over 74 days in a single (calm-to-moderate) regime; nothing here transfers to small caps, other markets, or stress periods, and the identity level effect~---about $-3$ points on this corpus~---may differ where firm priors are stronger or weaker.
